\section{Related Work}
\label{sec:related_work}

\subsection{Vision-Language-Action Models and Manipulation Safety}
Vision-Language-Action foundation models have redefined robot learning by transferring multimodal representations from web-scale pretraining to sensorimotor control \cite{openvla2024,pi02024}. Recent efforts, such as TacClaw-Infinity \cite{tacclaw2026infinity}, demonstrate that incorporating high-resolution tactile observations into flow-matching policies substantially boosts task success on multi-embodiment manipulation benchmarks. However, foundation policies predominantly output end-effector poses or joint trajectories via unconstrained regression or diffusion objectives. They lack formal mechanisms to guarantee that generated action trajectories respect nonlinear contact wrench manifolds or frictional slip boundaries. When physical contact properties deviate abruptly from pretraining distributions, generative policies are susceptible to constraint violations, highlighting the necessity of modular, real-time safety filters.

\subsection{Tactile Sensing and Incipient Slip Feedback}
High-resolution optical tactile sensors, including GelSight \cite{yuan2017gelsight}, DIGIT \cite{lambeta2020digit}, and GelSlim, capture dense geometric deformation and shear strain fields at the contact interface. Prior research has demonstrated that monitoring the divergence of marker displacement vectors provides an incipient slip detection window of $25\text{--}32\text{ ms}$ prior to macroscopic breakaway \cite{pandas2024incipient,shear2025grasp}. However, existing tactile controllers have largely coupled this pre-slip signal to heuristic proportional clamping or reactive thresholding reflexes. Such reactive reflexes act without forward predictive dynamics: by the time macroscopic slip is detected, the object has acquired substantial kinetic momentum, causing slip runaway under aggressive payloads. In contrast, our approach utilizes incipient slip to trigger predictive anticipatory force modulation within a model-based quadratic program.

\subsection{Control Barrier Functions for Contact Systems}
Control Barrier Functions (CBFs) provide mathematically rigorous forward-invariance certificates for safe sets \cite{ames2016control}. While continuous-time CBFs have been widely formulated for collision avoidance and bipedal walking, extending CBFs to multi-contact robotic grasping introduces substantial challenges. Contact forces do not appear directly as control inputs, requiring forward dynamics sensitivity mappings. Furthermore, sampled-data implementation introduces discretization error, rendering continuous-time invariance guarantees fragile unless discrete-time formulations are used \cite{agrawal2017discrete}. While prior work has explored continuous CBFs for two-jaw parallel grippers with rigid contact, this paper establishes a discrete-time predictive p-CBF framework for multi-fingered anthropomorphic hands interacting with compliant, dynamic payloads.

\subsection{Conformal Prediction and Non-Exchangeable Calibration}
Conformal prediction provides a distribution-free, finite-sample framework for quantifying predictive uncertainty. Standard split conformal methods require data exchangeability, an assumption routinely broken in closed-loop robotics due to state feedback and nonstationary environmental drift. Gibbs and Cand{\`e}s developed Adaptive Conformal Inference (ACI) \cite{gibbs2021adaptive,gibbs2024conformal}, establishing that online quantile adaptation guarantees long-run average coverage under arbitrary, non-exchangeable distribution shifts. While ACI has been explored for trajectory forecasting, this work represents the first integration of ACI with discrete-time Control Barrier Functions for certified contact-force regulation.
